\documentclass{exam}

\newif\ifA
\newif\ifkey
\newif\ifprintselected   % required by exam_config's \ansbox

%%%%%%%%%%%%%%%%%%%%%%%%%%%%%%%%%%%%%%%%%%%%%%%%%%%%%%
%%% SET THESE IF VARIABLES
%%%%%%%%%%%%%%%%%%%%%%%%%%%%%%%%%%%%%%%%%%%%%%%%%%%%%%
% Lab name -- labs are named by the lesson they belong to, never by a
% sequential lab number (Neil, 2026-10-12): this is the L6 Lab.
\def\labLesson{6}

% is this the key or not?
\ifdefined\iskey
	\ifnum\iskey=1
		\keytrue
	\else
		\keyfalse
	\fi
\else
	%% Pick one manually
	\keyfalse % if this is NOT the key
%	\keytrue     % if this IS the key
\fi

% set up the header and footer - change for every term
\ifdefined\thisTermIn
	\def\thisterm{\thisTermIn}
\else
	\def\thisterm{Spring 2027}
\fi
%
\def\thisclass{ECE 448}

%%%%%%%%%%%%%%%%%%%%%%%%%%%%%%%%%%%%%%%%%%%%%%%%%%%%%%
%%% END SET THESE IF VARIABLES
%%%%%%%%%%%%%%%%%%%%%%%%%%%%%%%%%%%%%%%%%%%%%%%%%%%%%%
% packages and shortcuts
\input{myPackages_gr}
\usepackage[hidelinks, linkcolor=red]{hyperref}
\input{myShortcuts}
\setlength{\parindent}{0pt}
\setlength{\parskip}{6pt}

%% exam class customization
\input{exam_config}
% ECE 448 document fonts.
%
% Body text is Barlow (SIL OFL; the faces and the license are in
% latex/fonts/Barlow/). Barlow is not in TeX Live, so it is loaded as an
% OpenType face through fontspec -- which means the document must be built
% with lualatex, not pdflatex. build_practice.sh and build_lab.sh both do.
%
% Compiling with pdflatex anyway is not an error: the guard below falls back
% to the LaTeX default so the source still compiles under any engine and can
% be syntax-checked. The resulting PDF just is not the one we ship.
%
% Paths are relative to latex/, which is where both build scripts run
% (each does cd "$(dirname "$0")"). Loading by path rather than by installed
% family name means no system font install and no fontconfig, so a fresh
% clone builds a byte-comparable PDF anywhere.
%
% MATH IS DELIBERATELY LEFT IN COMPUTER MODERN. Barlow has no math companion,
% and CM math under sans body text is already the house pairing: the decks set
% words in Source Sans Pro and let MathJax render symbols in its CM-derived
% face. Moving math to a sans as well would mean unicode-math plus a sans math
% font (Fira Math being the plausible one) and re-checking every equation in
% the course for symbol coverage. That is a separate decision, not a side
% effect of this one.

\usepackage{iftex}
\ifLuaTeX
	\usepackage{fontspec}
	\defaultfontfeatures{Ligatures=TeX}
	\setmainfont{Barlow}[
		Path           = fonts/Barlow/,
		Extension      = .ttf,
		UprightFont    = *-Regular,
		ItalicFont     = *-Italic,
		BoldFont       = *-Bold,
		BoldItalicFont = *-BoldItalic,
	]
	\setsansfont{Barlow}[
		Path           = fonts/Barlow/,
		Extension      = .ttf,
		UprightFont    = *-Regular,
		ItalicFont     = *-Italic,
		BoldFont       = *-SemiBold,
		BoldItalicFont = *-BoldItalic,
	]
	% microtype earns its place here specifically because of Barlow. Its
	% glyphs are narrower than Computer Modern's, so justified prose sets
	% tighter and TeX has to squeeze interword space to fit; two paragraphs
	% in the L02/L03 sets went overfull on the first Barlow build. Protrusion
	% and font expansion absorb that -- both are lualatex features, so this
	% belongs inside the guard with everything else.
	\usepackage{microtype}

	% Mono stays Latin Modern: Barlow ships no monospaced face, and verbatim
	% and code need one that actually is fixed-pitch.
\fi
   % Barlow body text (needs lualatex)
\printselectedfalse      % full worked solutions in the key, not answers-only
\CorrectChoiceEmphasis{\color{red}}
\SolutionEmphasis{\color{red}}
\renewcommand\questionlabel{\textbf{\thequestion.}}
\renewcommand{\thepartno}{\alph{partno}}
\renewcommand{\partlabel}{(\thepartno)}

% Fillable table cell: blank in the student copy, the expected value in red on
% the KEY. Fixed-width centered columns keep the blank cells writable.
\newcolumntype{K}[1]{>{\centering\arraybackslash}p{#1}}
\newcommand{\kv}[1]{\ifkey{\color{red}#1}\fi\rule[-0.10in]{0pt}{0.34in}}
% A shorter cell for a settings table, where the left columns hold text.
\newcommand{\kvs}[1]{\ifkey{\color{red}#1}\fi\rule[-0.08in]{0pt}{0.26in}}

% Answers are not in this public file. Every solution, \kv value, and \ansbox
% value is \keysol{<id>}, defined in the private ECE448-faculty repo at
% lab-keys/src/ECE448_Lab_L06_AMSimulation_solutions.tex; build_lab.sh puts that
% folder on TEXINPUTS (KEYDIR) for the KEY. A KEY build without it stops.
\long\def\defkeysol#1#2{\expandafter\long\expandafter\def\csname keysol@#1\endcsname{#2}}
\def\keysol#1{\ifcsname keysol@#1\endcsname\csname keysol@#1\endcsname\fi}
\ifkey
	\IfFileExists{ECE448_Lab_L06_AMSimulation_solutions.tex}%
		{\input{ECE448_Lab_L06_AMSimulation_solutions}}%
		{\errmessage{KEY build needs ECE448_Lab_L06_AMSimulation_solutions.tex from ECE448-faculty/lab-keys/src (set KEYDIR)}}
\fi

% print the answers if this is the key
\ifkey
	\printanswers % if this IS the key
\else
	\noprintanswers  %if this is NOT the key
\fi

% print the header and footer
% need to print key on top if this is the key
\ifkey
	\firstpageheader{}{\color{red}{ KEY}}{}
	\runningheader{\thisclass\hspace{1pt} -- \thisterm }{L\labLesson\ Lab\hspace{1pt}\color{red}{ KEY}}{\thepage\hspace{1pt} of\hspace{1pt} \numpages}
	\firstpagefooter{}{}{}
	\runningfooter{}{}{}
\else
	\firstpageheader{}{}{}
	\runningheader{\thisclass \hspace{1pt} -- \thisterm }{L\labLesson\ Lab\hspace{1pt}}{\thepage\hspace{1pt} of\hspace{1pt} \numpages}
	\firstpagefooter{}{}{}
	\runningfooter{}{}{}
\fi

\runningheadrule


\begin{document}
\pagestyle{headandfoot}   % exam class

%%%%%%%%%%%%%%%%%%%%%%%%%%%%%%%%%%%%%%%%%%%%%%%%%%%%%%
% title block
%%%%%%%%%%%%%%%%%%%%%%%%%%%%%%%%%%%%%%%%%%%%%%%%%%%%%%

\begin{center}
USAF ACADEMY \\
DEPARTMENT OF ELECTRICAL AND COMPUTER ENGINEERING \\
LESSON \labLesson\ LAB - AM Simulation \\
\textbf{Lab Packet} \\
\thisterm \\[4mm]
\end{center}

\vspace{-4pt}
\textbf{Name:} \rule[-2pt]{2.6in}{0.4pt} \hfill
\textbf{Section:} \rule[-2pt]{0.9in}{0.4pt} \hfill
\textbf{Date:} \rule[-2pt]{1.1in}{0.4pt}

\vspace{4pt}

\fullwidth{\textbf{Learning Objective 3.1: } I can describe DSB-AM, DSB-SC, and SSB
and identify each on a spectrum.
\\[2pt]\textbf{Learning Objective 3.2: } I can implement an AM modulator and
demodulator in GNU Radio and verify demodulated audio against the source.}

\vspace{2pt}

This packet is the L\labLesson\ Lab, in three parts, and you turn it in once, complete,
before Lesson 7. Do \textbf{Part 1} on your own before the lab period: list the
spectral lines you expect and the block settings you will need. \textbf{Part 2} is
your own measurements, on your own laptop, recorded in class: turn the sliders and
record what you see and hear. Do \textbf{Part 3} on your own after the lab: explain
what changed and why, in words.

\textbf{Grading.} One submission, graded against Objectives 3.1 and 3.2 on the
course's 1 / 0 scale, with revision allowed under course policy. Paste screenshots
where a question asks for one.

\textbf{No hardware.} Every flowgraph in this lab is a simulation. Nothing transmits,
and you need only your laptop and its speakers or headphones. Start with the volume
low.

The flowgraphs are on the L6 lesson page: \texttt{L06-dsb-am.grc},
\texttt{L06-dsb-sc.grc}, and \texttt{L06-ssb.grc}. They run the Lesson 6 message on a
carrier at $f_c = 10$ kHz, sampled at 48 kHz. \textbf{This lab uses its own message}:
two cosine tones of equal amplitude, at 1.5 kHz and 2.25 kHz, scaled so that their
sum peaks at exactly 1, on the same 10 kHz carrier. You set it in your own build and
in each downloaded flowgraph.

\begin{questions}

\fullwidth{\large\textbf{Part 1: Pre-Lab} \normalsize(individual, before the lab period)}

%%%%%%%%%%%%%%%%%%%%%%%%%%%%%%%%%%%%%%%%%%%%%%%%%%%
\question \textbf{Getting ready.} Work these for the lab's message before you build.

\begin{parts}

	\part List the frequency of every line you expect on the DSB-AM spectrum, and
	say which one is the carrier.

	\begin{center}
	\renewcommand{\arraystretch}{1.25}
	\begin{tabular}{| l | K{0.68in} | K{0.68in} | K{0.68in} | K{0.68in} | K{0.68in} |}
	\hline
	Line (kHz) & \kv{\keysol{kv1}} & \kv{\keysol{kv2}} & \kv{\keysol{kv3}} & \kv{\keysol{kv4}} & \kv{\keysol{kv5}} \\\hline
	\end{tabular}
	\end{center}

		\begin{solution}[0.6in]\keysol{sol1}\end{solution}

	\part List the lines you expect for DSB-SC, USB, and LSB.
		\begin{solution}[0.7in]\keysol{sol2}\end{solution}

	\newpage
	\part Fill in the settings you will type into the two message Signal Sources, in
	your build and in each downloaded flowgraph.

	\begin{center}
	\renewcommand{\arraystretch}{1.25}
	\begin{tabular}{| l | l | K{1.6in} |}
	\hline
	Block & Setting & Value \\\hline
	Message Signal Source 1 & Frequency (Hz) & \kvs{\keysol{kv6}} \\\hline
	Message Signal Source 1 & Amplitude & \kvs{\keysol{kv7}} \\\hline
	Message Signal Source 2 & Frequency (Hz) & \kvs{\keysol{kv8}} \\\hline
	Message Signal Source 2 & Amplitude & \kvs{\keysol{kv9}} \\\hline
	\end{tabular}
	\end{center}

		\begin{solution}[0.6in]\keysol{sol3}\end{solution}

\end{parts}

%%%%%%%%%%%%%%%%%%%%%%%%%%%%%%%%%%%%%%%%%%%%%%%%%%%
\newpage
\fullwidth{\large\textbf{Part 2: Measurements} \normalsize(your own laptop, in class)}

\question \textbf{Build the DSB-AM flowgraph.} Start GRC with an empty canvas. Do not
open \texttt{L06-dsb-am.grc}: build the transmitter and the envelope detector
yourself, with a \texttt{samp\_rate} of 48000, an \texttt{fc} of 10000, the message
settings from Question 1(c), a QT GUI Range named \texttt{m} from 0 to 1.5, a frequency
sink, a time sink comparing the detector
output with the delayed envelope, and an Audio Sink.

\begin{parts}

	\part Paste a screenshot of your flowgraph on the GRC canvas. Every connection
	must be black, not red.
		\begin{solution}[1.6in]\keysol{sol4}\end{solution}

	\part Run it and confirm the lines you listed in Question 1(a). Then set $m$ to each
	value below and read how many dB the 1.5 kHz tone's upper-sideband line sits below the
	carrier on the frequency sink.

	\begin{center}
	\renewcommand{\arraystretch}{1.25}
	\begin{tabular}{| l | K{0.9in} | K{0.9in} | K{0.9in} |}
	\hline
	$m$ & $0.25$ & $0.5$ & $1.0$ \\\hline
	Line below carrier (dB) & \kv{\keysol{kv10}} & \kv{\keysol{kv11}} & \kv{\keysol{kv12}} \\\hline
	\end{tabular}
	\end{center}

		\begin{solution}[0.5in]\keysol{sol5}\end{solution}

	\part Paste a screenshot of the time sink that compares the detector output with
	the delayed envelope at $m = 0.5$, and state whether they match.
		\begin{solution}[1.4in]\keysol{sol6}\end{solution}

\end{parts}

%%%%%%%%%%%%%%%%%%%%%%%%%%%%%%%%%%%%%%%%%%%%%%%%%%%
\newpage
\question \textbf{Overmodulation.} Keep your DSB-AM flowgraph running with the audio on.

\begin{parts}

	\begin{minipage}[t]{0.58\linewidth}
		\part Raise $m$ in steps of 0.05 and record the first value at which the detector
		output visibly departs from the envelope.
		\begin{solution}[0.7in]\keysol{sol7}\end{solution}
	\end{minipage}\hfill%
	\begin{minipage}[t]{0.37\linewidth}
		\raggedleft \vspace{12pt}
		$m = $ \ansbox[1.0in]{\keysol{ans1}}
	\end{minipage}

	\part Paste a screenshot at $m = 1.4$ and describe how the audio changes between
	$m = 0.5$ and $m = 1.4$.
		\begin{solution}[1.6in]\keysol{sol8}\end{solution}

\end{parts}

%%%%%%%%%%%%%%%%%%%%%%%%%%%%%%%%%%%%%%%%%%%%%%%%%%%
\question \textbf{DSB-SC.} Open \texttt{L06-dsb-sc.grc}, set its two message Signal
Sources to your Question 1(c) settings, and run it.

\begin{parts}

	\part Confirm the lines you listed in Question 1(b), and state the level of the
	frequency sink at the carrier frequency.
		\begin{solution}[0.6in]\keysol{sol9}\end{solution}

	\part With the LO phase at $0^\circ$ the coherent receiver's peak is 1.0. Set the LO
	phase to each value below, read the peak off the time sink, and listen to the audio
	at each.

	\begin{center}
	\renewcommand{\arraystretch}{1.25}
	\begin{tabular}{| l | K{0.68in} | K{0.68in} | K{0.68in} | K{0.68in} |}
	\hline
	LO phase & $30^\circ$ & $45^\circ$ & $60^\circ$ & $75^\circ$ \\\hline
	Peak output & \kv{\keysol{kv13}} & \kv{\keysol{kv14}} & \kv{\keysol{kv15}} & \kv{\keysol{kv16}} \\\hline
	\end{tabular}
	\end{center}

		\begin{solution}[0.4in]\keysol{sol10}\end{solution}

	\part Switch Listen To from Coherent to Envelope, and describe the difference.
		\begin{solution}[0.9in]\keysol{sol11}\end{solution}

\end{parts}

%%%%%%%%%%%%%%%%%%%%%%%%%%%%%%%%%%%%%%%%%%%%%%%%%%%
\newpage
\question \textbf{SSB.} Open \texttt{L06-ssb.grc}, set its two message Signal Sources
to your Question 1(c) settings, and run it.

\begin{parts}

	\part For each setting of the Sideband menu, record which lines remain and how far
	below them the suppressed sideband sits.

	\begin{center}
	\renewcommand{\arraystretch}{1.25}
	\begin{tabular}{| l | K{1.6in} | K{1.6in} |}
	\hline
	Sideband & Lines (kHz) & Suppression (dB) \\\hline
	USB & \kv{\keysol{kv17}} & \kv{\keysol{kv18}} \\\hline
	LSB & \kv{\keysol{kv19}} & \kv{\keysol{kv20}} \\\hline
	\end{tabular}
	\end{center}

		\begin{solution}[0.5in]\keysol{sol12}\end{solution}

	\part With LSB selected, set the tuning error to $+120$ Hz. Paste the time sink and
	describe the audio.
		\begin{solution}[1.4in]\keysol{sol13}\end{solution}

	\part Find the largest tuning error, either sign, at which you cannot hear any
	pitch change from the original. This is the threshold of a noticeable shift, not
	of intelligibility.
		\begin{solution}[0.5in]\keysol{sol14}\end{solution}

\end{parts}

%%%%%%%%%%%%%%%%%%%%%%%%%%%%%%%%%%%%%%%%%%%%%%%%%%%
\newpage
\fullwidth{\large\textbf{Part 3: Writeup} \normalsize(individual, after the lab)}

\question \textbf{Analysis.}

\begin{parts}

	\part Using your table in Question 2(b), describe what raising $m$ does to the
	sideband lines and to the carrier line, and explain in words why the carrier does
	not change.
		\begin{solution}[1.0in]\keysol{sol15}\end{solution}

	\part Explain in words why an envelope detector recovers DSB-AM but not DSB-SC, and
	what the coherent receiver has that the envelope detector does not.
		\begin{solution}[1.3in]\keysol{sol16}\end{solution}

	\part A DSB-SC LO that is off in phase loses volume, while an SSB receiver that is
	off in frequency changes pitch. Using your measurements, explain which error a
	voice listener tolerates better, and why SSB receivers have a tuning knob instead
	of carrier recovery.
		\begin{solution}[1.3in]\keysol{sol17}\end{solution}

\end{parts}

\end{questions}

\vspace{1.0em}
\noindent\textbf{Documentation:} \rule[-2pt]{0.6\linewidth}{0.4pt}

\end{document}
